\section{例子与应用}

\subsection{例 1：对公平硬币下注}

我们必须就一次公平掷硬币的结果押正面或反面；猜对赢一美元，猜错则输。假设我们押正面并赢了一美元，且未曾瞥见硬币的实际结果。我们的下注是赢钱的必要原因（或充分原因，或二者兼是）吗？

此例与 \S 1.4.4 讨论的临床试验（图 1.6）同构。以 $x$ 表示“我们押正面”，$y$ 表示“我们赢一美元”，$u$ 表示“硬币掷出正面”。$y$、$x$ 与 $u$ 之间的函数关系为
\begin{equation}
y=(x\wedge u)\vee(x'\wedge u'), \tag{9.36}
\end{equation}
它不是单调的，但由于模型被完全设定，我们可以由定义 \textup{(9.1)}—\textup{(9.3)} 计算因果概率。举例来说，
\[
\mathrm{PN}=P(y'_{x'}\mid x,y)=P(y'_{x'}\mid u)=1,
\]
因为 $x\wedge y\Rightarrow u$ 且 $Y_{x'}(u)=\text{false}$。用文字说：知道当前的下注（$x$）与当前的输赢（$y$），使我们能推断硬币的结果必为正面（$u$），由此又可推出：若押反面（$x'$）而非正面，本会输钱。同样，
\[
\mathrm{PS}=P(y_x\mid x',y')=P(y_x\mid u)=1
\]
（因为 $x'\wedge y'\Rightarrow u$），且
\begin{align*}
\mathrm{PNS} &= P(y_x,y'_{x'})\\
&= P(y_x,y'_{x'}\mid u)P(u)+P(y_x,y'_{x'}\mid u')P(u')\\
&= 1(0.5)+0(0.5)=0.5.
\end{align*}

我们看到，押正面有 50\% 的机会成为赢钱的必要且充分原因。然而，一旦我们赢了，便可百分之百地确定下注是赢钱的必要条件；一旦我们输了（比如押了反面），便可百分之百地确定押正面本会足以产生赢钱。这类反事实的经验内容见 \S 7.2.2 的讨论。

容易验证：若不知道 \textup{(9.36)} 中的函数关系——它告诉我们决定输赢的（确定性）政策——这些反事实量便无法由 $X$ 与 $Y$ 的联合概率算出（\S 1.4.4）。例如，从与本例相关的条件概率与因果效应
\[
P(y\mid x)=P(y\mid x')=P(y_x)=P(y_{x'})=P(y)=\tfrac12
\]
即可看出，因为同样的概率会由一个随机赔付政策产生，其中 $y$ 在函数上独立于 $x$——比如由一位盯着硬币、无视我们下注的庄家产生。在这样的随机政策中，因果概率 PN、PS 与 PNS 全为零。于是，按我们对可识别性的定义（定义 3.2.3），若两个模型在 $P$ 上一致而在某量 $Q$ 上不一致，则 $Q$ 不可识别。确实，定理 9.2.10 所刻画的界（式 \textup{(9.9)}）读作 $0\le\mathrm{PNS}\le\frac12$，意味着这三个因果概率无法仅由 $X$ 与 $Y$ 的统计数据确定，即便在受控实验中也不行；需要关于函数机制的知识，如 \textup{(9.36)} 所示。

有趣的是，硬币是在下注之前还是之后掷出，对刚才定义的因果概率没有影响。这与某些概率因果论（如 Good 1961）形成对照：后者试图通过把所有概率条件于所论原因（$x$）发生“之前一刻的世界状态”，来避开确定性机制。用于我们的下注故事时，其意图是把一切概率条件于硬币的状态（$u$），但若硬币在下注之后才掷出，这一意图便落空了。试图用原因之后发生的事件来丰富条件集的努力，最终又回到了涉及反事实变量的确定性关系（见 Cartwright 1989、Eells 1991，以及 \S 7.5.4 的讨论）。

当然，有人会争辩：若硬币在下注之后掷出，那么我们若改押别的、赢钱会怎样，就完全不清楚了；仅仅是说出我们的下注，就可能影响硬币的轨迹（Dawid 2000）。这一反驳可通过把 $x$ 与 $u$ 置于两处遥远的地点来化解：在下注后的一瞬间掷硬币，但要在任何光线来得及从下注室传到掷币室之前。在这样的假想情境中，“若我们改押，赢钱会不同”这一反事实陈述颇有说服力，尽管条件事件（$u$）发生在所论原因（$x$）之后。我们由此得出结论：“$x$ 之前一刻的世界状态”这类时间性描述，不能用来恰当地识别问题中合适的条件事件集（$u$）；要表述“因果概率”这一概念，需要关于所涉机制的确定性模型。

\subsection{例 2：行刑队}

再考虑 \S 7.1.2 的行刑队（见图 9.1）；$A$ 与 $B$ 是步枪手，$C$ 是队长（他正在等待法庭命令 $U$），$T$ 是死囚。设 $u$ 为“法庭已下令处决”这一命题，$x$ 为“$A$ 扣动扳机”这一命题，$y$ 为“$T$ 已死”。我们再次假定 $P(u)=\frac12$，$A$ 与 $B$ 是警觉且守法的神枪手，且 $T$ 不大可能因惊吓或其他外因而死。我们希望计算 $x$ 是 $y$ 之必要（或充分，或二者兼是）原因的概率（即计算 PN、PS 与 PNS）。

\begin{figure}[htbp]
\centering
\begin{tikzpicture}[causal,scale=1.0,every node/.style={transform shape}]
  \node[vtxs] (U) at (0,3.0) {$U$};
  \node[vtxs] (C) at (0,2.1) {$C$};
  \node[vtxs] (A) at (-1.5,0.9) {$A$};
  \node[vtxs] (B) at (1.5,0.9) {$B$};
  \node[vtxs] (T) at (0,-0.3) {$T$};
  \draw[ed] (U)--(C);
  \draw[ed] (C)--(A); \draw[ed] (C)--(B);
  \draw[ed] (A)--(T); \draw[ed] (B)--(T);
  \node[font=\footnotesize,anchor=east] at (-2.05,0.9) {$x$：$A$ 开枪};
  \node[font=\footnotesize,anchor=west] at (0.5,3.0) {（法庭命令）};
  \node[font=\footnotesize,anchor=west] at (0.5,2.1) {（队长）};
  \node[font=\footnotesize,anchor=west] at (1.85,0.9) {（步枪手）};
  \node[font=\footnotesize,anchor=west] at (0.5,-0.3) {$y$：$T$ 死亡};
\end{tikzpicture}
\caption{两人行刑队例子中的因果关系。}
\label{fig:9.1}
\end{figure}

定义 9.2.1—9.2.3 使我们能直接从给定的因果模型计算这些概率，因为一切函数与一切概率都已指定，每个变量的真值都随 $U$ 而定。据此可写\footnote{回忆 $P(Y_x(u')=\text{true})$ 涉及子模型 $M_x$，其中 $X$ 被独立于 $U$ 地设为“真”。因此，尽管在条件 $u'$ 下队长并未发出信号，潜在结局仍要求我们假设步枪手 $A$ 违法扣动了扳机（$x$）。}
\begin{equation}
\begin{aligned}
P(y_x) &= P(Y_x(u)=\text{true})P(u)+P(Y_x(u')=\text{true})P(u')\\
&= \tfrac12(1+1)=1.
\end{aligned} \tag{9.37}
\end{equation}
同样，
\begin{equation}
\begin{aligned}
P(y_{x'}) &= P(Y_{x'}(u)=\text{true})P(u)+P(Y_{x'}(u')=\text{true})P(u')\\
&= \tfrac12(1+0)=\tfrac12.
\end{aligned} \tag{9.38}
\end{equation}
为计算 PNS，我们必须求联合事件 $y_x\wedge y'_{x'}$ 的概率。由于这两个事件只在 $U=\text{true}$ 时共同为真，我们有
\begin{equation}
\begin{aligned}
\mathrm{PNS} &= P(y_x,y'_{x'})\\
&= P(y_x,y'_{x'}\mid u)P(u)+P(y_x,y'_{x'}\mid u')P(u')\\
&= \tfrac12(0+1)=\tfrac12.
\end{aligned} \tag{9.39}
\end{equation}

PS 与 PN 的计算同样因下述事实而简化：每个条件事件——PN 的 $x\wedge y$、PS 的 $x'\wedge y'$——只在 $U$ 的一个状态下为真。于是有
\[
\mathrm{PN}=P(y'_{x'}\mid x,y)=P(y'_{x'}\mid u)=0,
\]
这反映出：一旦法庭下令处决（$u$），$T$ 会因步枪手 $B$ 的射击而死（$y$），即便 $A$ 未开枪（$x'$）。的确，得知 $T$ 之死后，我们可以断然指出：步枪手 $A$ 的射击不是死亡的必然原因。同样，
\[
\mathrm{PS}=P(y_x\mid x',y')=P(y_x\mid u')=1,
\]
与我们的直觉相符：专家射手的一枪足以导致 $T$ 死亡，无论法庭如何裁决。

注意，定理 9.2.10 与 9.2.11 不适用于此例，因为 $x$ 不是外生的；事件 $x$ 与 $y$ 有共同原因（队长的信号），这使得 $P(y\mid x')=0\ne P(y_{x'})=\frac12$。然而，$Y$（关于 $x$）的单调性使我们能由联合分布 $P(x,y)$ 与因果效应（用 \textup{(9.28)}—\textup{(9.30)}）计算 PNS、PS 与 PN，而不必求助于函数模型。确实，写出
\begin{equation}
P(x,y)=P(x',y')=\tfrac12 \tag{9.40}
\end{equation}
与
\begin{equation}
P(x,y')=P(x',y)=0, \tag{9.41}
\end{equation}
我们得到
\begin{equation}
\mathrm{PN}=\frac{P(y)-P(y_{x'})}{P(x,y)}=\frac{\tfrac12-\tfrac12}{\tfrac12}=0 \tag{9.42}
\end{equation}
与
\begin{equation}
\mathrm{PS}=\frac{P(y_x)-P(y)}{P(x',y')}=\frac{1-\tfrac12}{\tfrac12}=1, \tag{9.43}
\end{equation}
正如所料。

\subsection{例 3：辐射对白血病的影响}

考虑下述数据（表 9.1，改编自 Finkelstein and Levin 1990\footnote{Finkelstein and Levin (1990) 中的数据以“人年”为单位。为作说明，我们假定十年观察期，把数据换算为绝对数（死亡与未死亡人数）。}），它们比较了犹他州南部儿童在高暴露与低暴露于内华达核试验沉降物辐射之下的白血病死亡率。给定这些数据，我们希望估计高暴露于辐射是白血病死亡之必要（或充分，或二者兼是）原因的概率。

\begin{table}[htbp]
\centering
\caption{高暴露（$x$）与低暴露（$x'$）儿童的白血病死亡人数}
\label{tab:9.1}
\small
\begin{tabular}{lcc}
\toprule
& \multicolumn{2}{c}{暴露}\\
\cmidrule(lr){2-3}
& 高（$x$） & 低（$x'$）\\
\midrule
死亡（$y$）   & 30     & 16\\
存活（$y'$）  & 69,130 & 59,010\\
\bottomrule
\end{tabular}
\end{table}

假定单调性——即暴露于核辐射对研究中的任何个体都没有治疗作用——该过程可由一个简单的析取机制建模，其方程为
\begin{equation}
y=f(x,u,q)=(x\wedge q)\vee u, \tag{9.44}
\end{equation}
其中 $u$ 表示 $y$ 的“一切其他原因”，$q$ 表示为使 $x$ 触发 $y$ 所必须具备的一切“促成”机制。假定 $q$ 与 $u$ 均不可观测，我们要问：在什么条件下可以由 $X$ 与 $Y$ 的联合分布识别因果概率（PNS、PN 与 PS）。

由于 \textup{(9.44)} 关于 $x$ 单调，定理 9.2.14 指出：只要 $X$ 是外生的，三个量便都可识别；即 $x$ 应与 $q$、$u$ 独立。在此假设下，\textup{(9.21)}—\textup{(9.23)} 进一步使我们能由频率数据计算因果概率。以分数代表概率，表 9.1 的数据给出如下数值结果：
\begin{equation}
\mathrm{PNS}=P(y\mid x)-P(y\mid x')=\frac{30}{30+69{,}130}-\frac{16}{16+59{,}010}=0.0001625, \tag{9.45}
\end{equation}
\begin{equation}
\mathrm{PN}=\frac{\mathrm{PNS}}{P(y\mid x)}=\frac{\mathrm{PNS}}{30/(30+69{,}130)}=0.37535, \tag{9.46}
\end{equation}
\begin{equation}
\mathrm{PS}=\frac{\mathrm{PNS}}{1-P(y\mid x')}=\frac{\mathrm{PNS}}{1-16/(16+59{,}010)}=0.0001625. \tag{9.47}
\end{equation}

从统计学上说，这些数字意味着：

\begin{enumerate}[label=\arabic*.,leftmargin=1.8em]
\item 随机抽取的一名儿童，若暴露则死于白血病、若不暴露则存活的概率为万分之 1.625；
\item 一名死于白血病的暴露儿童，若未曾暴露本会存活的概率为 37.544\%；
\item 任何未暴露且存活的儿童，若曾暴露本会死于白血病的概率为万分之 1.625。
\end{enumerate}

Glymour (1998) 分析此例，意在识别概率 $P(q)$（Cheng 的“因果力”），它与 PS 一致（见引理 9.2.8）。Glymour 的结论是：只要 $x$、$u$ 与 $q$ 相互独立，$P(q)$ 便可识别并由 \textup{(9.23)} 给出。我们的分析表明，Glymour 的结果可在若干方面推广。第一，由于 $Y$ 关于 $X$ 单调，即便 $q$ 与 $u$ 相依，\textup{(9.23)} 的有效性仍有保证，因为外生性只要求 $x$ 与 $\{u,q\}$ 联合独立。这在流行病学情境中很重要，因为个体对核辐射的易感性很可能与对白血病其他潜在原因（如天然辐射）的易感性相关。

第二，定理 9.2.11 向我们保证：PN、PS 与 PNS 之间的关系（式 \textup{(9.11)}—\textup{(9.12)}）——Glymour 对独立的 $q$ 与 $u$ 导出它们——在 $u$ 与 $q$ 相依时仍应成立。

最后，定理 9.2.15 向我们保证：即便 $x$ 与 $\{u,q\}$ 不独立，只要 \textup{(9.44)} 的机制嵌入一个更大的因果结构、使 $P(y_x)$ 与 $P(y_{x'})$ 可识别，PN 与 PS 仍可识别。例如，假定怀疑暴露于核辐射（$x$）与地形和海拔相关，而后者也是决定宇宙辐射暴露的因素。反映这种考量的模型如图 9.2 所示，其中 $W$ 表示同时影响 $X$ 与 $U$ 的因素。校正 $X$ 对 $Y$ 之因果效应中可能的混杂偏倚，一个自然的办法是调整 $W$，即用标准调整公式（式 \textup{(3.19)}）计算 $P(y_x)$ 与 $P(y_{x'})$
\begin{equation}
P(y_x)=\sum_w P(y\mid x,w)P(w),\qquad P(y_{x'})=\sum_w P(y\mid x',w)P(w) \tag{9.48}
\end{equation}
（而非 $P(y\mid x)$ 与 $P(y\mid x')$），其中求和遍及 $W$ 的各个水平。这一调整公式由 \textup{(9.35)} 推出，无论中介 $X$ 与 $Y$ 的机制如何都正确，只要 $W$ 代表影响 $X$ 与 $Y$ 的\emph{全部}共同因素（见 \S 3.3.1）。

定理 9.2.15 指示我们把 \textup{(9.48)} 代入 \textup{(9.29)} 与 \textup{(9.30)} 以评估 PN 与 PS，并保证所得表达式构成 PN 与 PS 的一致估计。这种一致性由单调性假设与因果图的（假定）拓扑共同保证。

注意，\textup{(9.20)} 定义的单调性是 $x$ 与 $y$ 之间一切通路的全局性质。因果模型在这些通路上可以包含若干非单调机制而不影响 \textup{(9.20)} 的有效性。然而，支持单调性成立的论据必须基于实质信息，因为它一般不可检验。例如，Robins and Greenland (1989) 主张，暴露于核辐射可能对某些个体有益，因为此类辐射在临床上常规用于治疗癌症病人。\textup{(9.32)} 中的不等式构成单调性的一个统计检验（尽管是弱检验），它同时基于实验与观察研究。

\begin{figure}[htbp]
\centering
\begin{tikzpicture}[causal,scale=0.98,every node/.style={transform shape}]
  \node[vtx] (Q) at (-2.4,1.35) {$Q$};
  \node[vtx] (X) at (-0.2,1.15) {$X$};
  \node[vtx] (W) at (1.0,2.05) {$W$};
  \node[vtx] (U) at (1.9,-0.55) {$U$};
  \node[vtx] (Y) at (-0.9,-2.75) {$Y$};
  % AND gate
  \draw[semithick] (-2.1,0.3) -- (-2.1,-0.2) arc[start angle=180,end angle=360,x radius=1.2,y radius=0.55] -- (0.3,0.3) -- cycle;
  \node[font=\small] at (-0.9,-0.25) {AND};
  % OR gate
  \draw[semithick] (-1.55,-1.35) -- (-0.25,-1.35) -- (-0.9,-2.1) -- cycle;
  \node[font=\small] at (-0.9,-1.62) {OR};
  % edges
  \draw[ed] (Q)--(-1.75,0.3);
  \draw[ed] (X)--(-0.2,0.3);
  \draw[ed] (W)--(X);
  \draw[ed] (W)--(U);
  \draw[ed] (-0.9,-0.75)--(-0.9,-1.35);
  \draw[ed] (U)--(-0.3,-1.35);
  \draw[ed] (-0.9,-2.1)--(Y);
  % latent association between Q and U
  \draw[<->,dashed,semithick] (Q) to[out=95,in=55] (U);
  % labels
  \node[font=\footnotesize,anchor=east] at (-2.75,1.35) {促成因素};
  \node[font=\footnotesize,anchor=west] at (1.65,2.25) {混杂因子};
  \node[font=\footnotesize,anchor=west] at (2.25,-0.55) {其他原因};
  \node[font=\footnotesize,anchor=east] at (-0.62,1.12) {辐射};
  \node[font=\footnotesize,anchor=west] at (-0.5,-2.75) {白血病};
\end{tikzpicture}
\caption{辐射—白血病例子中的因果关系，其中 $W$ 表示混杂因子。}
\label{fig:9.2}
\end{figure}

\subsection{例 4：由实验与非实验数据确定法律责任}

一桩诉讼针对药物 $x$ 的制造商提起，指控该药很可能导致了 A 先生的死亡；A 先生服用该药以缓解与疾病 $D$ 相关的症状 $S$。

制造商声称，关于有症状 $S$ 之病人的实验数据确凿地表明，药物 $x$ 只可能使死亡率略微上升。然而原告辩称，该实验研究与此案关系不大，因为它代表的是药物对\emph{全部}病人的效应，而不是对像 A 先生这样在服用药物 $x$ 期间确实死亡的病人的效应。此外原告主张，A 先生有其独特之处：他出于自愿服药，不同于为遵守实验方案而服药的受试者。为支持这一论证，原告提供非实验数据，表明大多数选择药物 $x$ 的病人若不是因为该药，本会活着。制造商反驳说：(1) 关于病人本会或不会死亡的反事实推测纯属形而上学，应当避免（Dawid 2000）；(2) 非实验数据应被先验地摒弃，因为此类数据可能被无关因素严重混杂。法庭现在必须基于实验与非实验两项研究，判定药物 $x$ 事实上是 A 先生死因的概率。

与这两项研究相关的（假想）数据见表 9.2。实验数据给出估计
\begin{equation}
P(y_x)=16/1000=0.016, \tag{9.49}
\end{equation}
\begin{equation}
P(y_{x'})=14/1000=0.014; \tag{9.50}
\end{equation}
非实验数据给出估计
\begin{equation}
P(y)=30/2000=0.015, \tag{9.51}
\end{equation}
\begin{equation}
P(y,x)=2/2000=0.001. \tag{9.52}
\end{equation}
把这些估计代入 \textup{(9.29)}——它给出 PN 的下界（见 \textup{(11.42)}）——得到
\begin{equation}
\mathrm{PN}\ge\frac{P(y)-P(y_{x'})}{P(y,x)}=\frac{0.015-0.014}{0.001}=1.00. \tag{9.53}
\end{equation}
可见原告是对的；除抽样误差外，数据以 100\% 的把握告诉我们：药物 $x$ 事实上对 A 先生之死负有责任。注意，直接使用实验的超额风险比会得到一个低得多（且不正确）的结果：
\begin{equation}
\frac{P(y_x)-P(y_{x'})}{P(y_x)}=\frac{0.016-0.014}{0.016}=0.125. \tag{9.54}
\end{equation}

\begin{table}[htbp]
\centering
\caption{实验研究与非实验研究的数据}
\label{tab:9.2}
\small
\begin{tabular}{lcccc}
\toprule
& \multicolumn{2}{c}{实验} & \multicolumn{2}{c}{非实验}\\
\cmidrule(lr){2-3}\cmidrule(l){4-5}
& $x$ & $x'$ & $x$ & $x'$\\
\midrule
死亡（$y$）  & 16  & 14  & 2   & 28\\
存活（$y'$） & 984 & 986 & 998 & 972\\
\bottomrule
\end{tabular}
\end{table}

显然，实验研究未揭示的是：若有选择，临终病人会避开药物 $x$。的确，若存在任何（在有选择时）会选择 $x$ 的临终病人，那么对照组（$x'$）本会包含一些这样的病人（由于随机化），于是对照组的死亡比例本会高于 $P(x',y)$——即避开 $x$ 的临终病人所占的人群比例。然而，等式 $P(y_{x'})=P(y,x')$ 告诉我们，对照组中没有这样的病人；因此（由随机化）总体中也不存在这样的病人，从而自愿选择药物 $x$ 的病人中没有一个是临终病例；他们全都对 $x$ 易感。

表 9.2 的数字显然是为代表一个极端情形而刻意构造的，以便对 \textup{(9.29)} 的有效性作定性说明。尽管如此，值得注意的是：实验与非实验研究的结合可以揭示实验研究单独无法揭示的东西；此外，这种结合还可为实验程序是否恰当提供一项必要的检验。例如，若表 9.2 中的频率略有不同，它们很容易使 \textup{(9.53)} 中的 PN 取大于 1 的值，或违反 \textup{(9.33)} 的基本不等式。这种违反将表明实验组与非实验组不相容，或许源于抽样不当。

最后一点或许需要略作解释，以免读者疑惑：为什么两个取自不同实验条件下不同群体的数据集会互相约束？解释是：只要两个子群体都是从总体中适当抽样的，其中某些量就应当对这些差异全都保持不变。这些不变的量就是因果效应概率 $P(y_{x'})$ 与 $P(y_x)$。虽然这些反事实概率并未在观察组中测量，但（按定义）它们必定与实验组中测得者相同。这些量的不变性是受控实验的基本公理；没有它，从实验研究推断总体的一般行为便不可能。这些量的不变性蕴含 \textup{(9.33)} 的不等式；若单调性成立，则 \textup{(9.32)} 随之成立。

\subsection{结果小结}

现在小结 \S 9.2 与 \S 9.3 中对执业流行病学家与政策制定者有价值的结果。这些结果见表 9.3，它列出在各种假设与各类数据下 PN（对非实验事件）的最佳估计量——假设越强，估计越有信息量。

我们看到，流行病学家常等同于因果概率的超额风险比（ERR），只有在能够确认两个假设时才有效度量 PN：外生性（即无混杂）与单调性（即无预防）。当单调性不成立时，ERR 仅给出 PN 的下界，如 \textup{(9.13)} 所示。（上界通常是 1。）表 9.3 右端的空缺（——）代表空洞的界（即 $0\le\mathrm{PN}\le1$）。存在混杂时，ERR 必须加上附加项 $[P(y\mid x')-P(y_{x'})]/P(x,y)$ 加以修正，如 \textup{(9.31)} 所述。换言之，当（因果效应的）混杂偏倚为正时，PN 比 ERR 高出这一附加项。显然，由于除以 $P(x,y)$，PN 的偏倚可能比因果效应偏倚 $P(y\mid x')-P(y_{x'})$ 高出许多倍。然而混杂仅源于暴露与影响结果的其他因素之间的关联；人们无需关心这些因素与暴露易感性之间的关联（见图 9.2）。

\begin{table}[htbp]
\centering
\caption{PN 作为假设与可用数据的函数}
\label{tab:9.3}
\small
\begin{tabular}{lccccc}
\toprule
\multicolumn{3}{c}{假设} & \multicolumn{3}{c}{可用数据}\\
\cmidrule(lr){1-3}\cmidrule(l){4-6}
外生性 & 单调性 & 附加假设 & 实验 & 观察 & 合并\\
\midrule
$+$ & $+$ &            & ERR & ERR & ERR\\
$+$ & $-$ &            & 界  & 界  & 界\\
$-$ & $+$ & 协变量控制 & —— & 修正 ERR & 修正 ERR\\
$-$ & $+$ &            & —— & —— & 修正 ERR\\
$-$ & $-$ &            & —— & —— & 界\\
\bottomrule
\end{tabular}

\par\smallskip
{\footnotesize 注：ERR 表示超额风险比，即 $1-P(y\mid x')/P(y\mid x)$；修正 ERR 见 \textup{(9.31)}。}
\end{table}

表 9.3 的最后一行对应完全不作假设的情形，它导致 PN 的空洞界，除非我们有合并数据。然而这并不意味着，除单调性与外生性之外的其他正当假设无助于使 PN 可识别。此类假设的运用将在下一节探讨。

\section{非单调模型中的识别}

本节讨论不作单调性假设时因果概率的识别。我们将假定给定一个因果模型 $M$，其中一切函数关系均已知；但由于背景变量 $U$ 不可观测，其分布未知，模型设定并不完整。

我们的第一步是研究在什么条件下函数 $P(u)$ 可被识别，从而使整个模型可识别。若 $M$ 是马尔可夫的，则问题可通过分别考察每个父—子家族来分析。考虑 $M$ 中任一方程
\begin{equation}
y=f(\pa_Y,u_Y)=f(x_1,x_2,\dots,x_k,u_1,\dots,u_m), \tag{9.55}
\end{equation}
其中 $U_Y=\{U_1,\dots,U_m\}$ 是出现在 $Y$ 的方程中的背景（可能相依的）变量集。一般地，$U_Y$ 的定义域可以是任意的、离散的或连续的，因为这些变量代表被略去于模型之外的未观测因素。然而，由于可观测变量是二值的，从 $\PA_Y$ 到 $Y$ 的函数只有有限多个，且对任一点 $U_Y=u$，其中只有一个函数被实现。这定义了 $U_Y$ 定义域的一个典范划分，把它分成等价类之集 $S$，其中每个等价类 $s\in S$ 诱导同一个从 $\PA_Y$ 到 $Y$ 的函数 $f^{(s)}$（见 \S 8.2.2）。于是，当 $u$ 在其定义域上变动时，一个函数集 $S$ 被实现，我们可以把 $S$ 视为新的背景变量，其取值对应于在 $U_Y$ 中可实现的那个从 $\PA_Y$ 到 $Y$ 的函数集 $\{f^{(s)}:s\in S\}$。此类函数的数目通常小于 $2^{(2^k)}$。\footnote{Balke and Pearl (1994a,b) 把这些 $S$ 变量称为“响应变量”，如 \S 8.2.2 那样；Heckerman and Shachter (1995) 称之为“映射变量”。}

例如，考虑图 9.2 所描述的模型。当背景变量 $(Q,U)$ 在各自定义域上变动时，$X$ 与 $Y$ 之间的关系跨越三个不同的函数：
\[
f^{(1)}:\;Y=\text{true},\qquad f^{(2)}:\;Y=\text{false},\qquad f^{(3)}:\;Y=X.
\]
第四个可能的函数 $Y\ne X$ 永远不会被实现，因为 $f_Y(\cdot)$ 是单调的。单元格 $(q,u)$ 与 $(q',u)$ 诱导 $X$ 与 $Y$ 之间相同的函数，因此属于同一个等价类。

若给定分布 $P(u_Y)$，我们便可算出分布 $P(s)$，而后者将确定条件概率 $P(y\mid\pa_Y)$：把 $P(s)$ 对一切把 $\pa_Y$ 映射为 true 的函数 $f^{(s)}$ 求和，
\begin{equation}
P(y\mid\pa_Y)=\sum_{s:\,f^{(s)}(\pa_Y)=\text{true}}P(s). \tag{9.56}
\end{equation}
为确保模型可识别，充分条件是：我们能反转这一过程，由 $P(y\mid\pa_Y)$ 确定 $P(s)$。若把条件概率集 $P(y\mid\pa_Y)$ 表示为一个向量 $\vec{p}$（维数为 $2^k$），把 $P(s)$ 表示为向量 $\vec{q}$，则 \textup{(9.56)} 定义了 $\vec{p}$ 与 $\vec{q}$ 之间的线性关系，可表示为矩阵乘法（如 \textup{(8.13)} 那样）
\begin{equation}
\vec{p}=R\vec{q}, \tag{9.57}
\end{equation}
其中 $R$ 是一个 $2^k\times|S|$ 矩阵，其元素或为 0 或为 1。于是，识别的一个充分条件很简单：$R$ 连同归一化方程 $\sum_j q_j=1$ 可逆。

一般地，$R$ 不可逆，因为 $\vec{q}$ 的维数可能远大于 $\vec{p}$ 的维数。然而在许多情形中，例如“噪声或”机制
\begin{equation}
Y=U_0\;\bigvee_{i=1,\dots,k}(X_i\wedge U_i), \tag{9.58}
\end{equation}
其对称性使得即便外生变量 $U_0,U_1,\dots,U_k$ 不独立，$\vec{q}$ 也可由 $P(y\mid\pa_Y)$ 识别。这可通过注意以下事实看出：每个满足 $U_0=\text{false}$ 的点 $u$ 都定义一个唯一函数 $f^{(s)}$，因为若 $T$ 是使 $U_i$ 为真的指标 $i$ 之集，则 $\PA_Y$ 与 $Y$ 的关系变为
\begin{equation}
Y=U_0\;\bigvee_{i\in T}X_i, \tag{9.59}
\end{equation}
而对 $U_0=\text{false}$，此方程为每个 $T$ 定义一个不同的函数。所诱导的函数数目是 $2^k+1$，减去 1 以作归一化，恰为 $\PA_Y$ 的不同实现之数。此外，容易证明连接 $\vec{p}$ 与 $\vec{q}$ 的矩阵可逆。我们由此断言：在任何由噪声或机制构成的马尔可夫模型中，每个反事实语句的概率都可识别，无论各家族中的背景变量是否相互独立。当然，对噪声与机制或其任意组合（包括取反机制）也成立，只要每个家族由同一类型的机制构成。

为把这一结果推广到噪声或与噪声与之外的机制，我们注意：尽管此例中的 $f_Y(\cdot)$ 是单调的（关于每个 $X_i$），但确保可识别的是 $f_Y(\cdot)$ 的\emph{冗余性}而非其单调性。下面是一个单调但 $R$ 矩阵不可逆的例子：
\[
Y=(X_1\wedge U_1)\vee(X_2\wedge U_1)\vee(X_1\wedge X_2\wedge U_3).
\]
当 $U_3=\text{false}$ 时，此函数代表噪声或门；当 $U_3=\text{true}$ 且 $U_1=U_2=\text{false}$ 时，它变成噪声与门。所诱导的等价类有六个，需要五个独立方程来确定其概率；而数据 $P(y\mid\pa_Y)$ 只提供四个这样的方程。

相反，由下述函数支配的机制虽然非单调，却是可逆的：
\[
Y=\mathrm{XOR}(X_1,\mathrm{XOR}(U_2,\dots,\mathrm{XOR}(U_{k-1},\mathrm{XOR}(X_k,U_k)))),
\]
其中 $\mathrm{XOR}(\cdot)$ 表示异或。此方程只诱导两个从 $\PA_Y$ 到 $Y$ 的函数：
\[
Y=\begin{cases}
\mathrm{XOR}(X_1,\dots,X_k) & \text{若 } \mathrm{XOR}(U_1,\dots,U_k)=\text{false},\\
\neg\,\mathrm{XOR}(X_1,\dots,X_k) & \text{若 } \mathrm{XOR}(U_1,\dots,U_k)=\text{true}.
\end{cases}
\]
因此，识别所需的参数只有一个，即 $P[\mathrm{XOR}(U_1,\dots,U_k)=\text{true}]$；而单个条件概率（比如 \mbox{$P(y\mid x_1,\dots,x_k)$}）就足以把它算出。

我们用一个定理来总结这些考量。

\begin{definition}[局部可逆性]
若对每个变量 $V_i\in V$，$2^k+1$ 个方程
\begin{equation}
P(y\mid\pa_i)=\sum_{s:\,f^{(s)}(\pa_i)=\text{true}}q_i(s), \tag{9.60}
\end{equation}
\begin{equation}
\sum_s q_i(s)=1 \tag{9.61}
\end{equation}
关于 $q_i(s)$ 有唯一解，则称模型 $M$ 是局部可逆的；其中每个 $f_i^{(s)}(\pa_i)$ 对应于 $u_i$ 在等价类 $s$ 中所诱导的函数 $f_i(\pa_i,u_i)$。
\end{definition}

\begin{theorem}
给定马尔可夫模型 $M=\langle U,V,\{f_i\}\rangle$，其中函数 $\{f_i\}$ 已知而外生变量 $U$ 不可观测。若 $M$ 局部可逆，则每个反事实语句的概率都可由联合概率 $P(v)$ 识别。
\end{theorem}

\begin{proof}
若 \textup{(9.60)} 关于 $q_i(s)$ 有唯一解，我们便可把 $U$ 替换为 $S$，得到如下等价模型：
\[
M'=\langle S,V,\{f_i'\}\rangle,\qquad\text{其中 } f_i'=f_i^{(s)}(\pa_i).
\]
模型 $M'$ 连同 $q_i(s)$ 完全指定了一个概率因果模型 $\langle M',P(s)\rangle$（由马尔可夫性），由此按定义可导出反事实的概率。
\end{proof}

定理 9.4.2 为识别因果概率提供了一个充分条件，但它当然没有穷尽有助于实现识别的假设。在许多情形中，我们可能有理由对模型假设额外的结构——例如，进入每个家族的 $U$ 变量本身相互独立。此时对概率 $P(s)$ 施加了额外约束，即便 $S$ 的基数远超条件概率 $P(y\mid\pa_Y)$ 的数目，\textup{(9.60)} 也可能有解。

\section{结论}

本章阐明并分析了因果的必要成分与充分成分之间的相互作用。借助立足于结构模型语义的反事实解释，我们演示了计算反事实概率的简单技术如何可用于计算因果概率、判定识别问题、揭示因果概率可由统计数据估计的条件，以及为分析者与研究者例行（往往不自觉地）作出的假设设计检验。

在实践方面，我们为流行病学家与健康科学家提供了若干有用工具（部分总结于表 9.3）。本章表述并提请注意一些微妙的假设，在使用超额风险比之类统计度量代表归因风险或因果概率之类因果量之前，必须先确认这些假设（定理 9.2.14）。它表明如何把实验与非实验研究的数据结合起来，以产生任何单一研究都无法揭示的信息（定理 9.2.15 与 \S 9.3.4）。最后，它为常用的“无预防”假设，以及“临床试验是否代表其目标人群”这一常被追问的问题，提供了检验（式 \textup{(9.32)}）。

在概念方面，我们看到必要性概率（PN）与充分性概率（PS）都参与我们对因果的理解，且每个成分都有其逻辑与计算规则。尽管必要原因的反事实概念（即结果“若不是”该行动就不会发生）在法律情境（Robertson 1997）与日常话语中占主导地位，因果的充分性成分对因果思维仍有确定的影响。

充分性成分的重要性，可在必要成分休眠或被确保的例子中揭示。为什么我们认为划火柴（对一场火灾而言）是比氧气的存在更恰当的解释？用 PN 与 PS 的语言重述这个问题，我们注意到：既然两种解释对火灾都是必要的，各自都会得到为 1 的 PN。（事实上，若允许其他点火方式，氧气的 PN 还更高。）因此，必定是充分性成分赋予了火柴比氧气更强的解释力。若划火柴与氧气存在所关联的概率分别记为 $p_m$ 与 $p_o$，则这些解释所对应的 PS 度量为 $\mathrm{PS}(\text{火柴})=p_o$ 与 $\mathrm{PS}(\text{氧气})=p_m$；当 $p_o>p_m$ 时，显然偏向火柴。于是，一个被指派去解释火灾何以发生的机器人，在它的考量中别无选择，只能同时考虑 PN 与 PS。

PS 是否应进入刑法与侵权法的考量？我相信应当——Good (1993) 也如此认为——因为关注充分性意味着关注一个人行动的后果。划火柴的人本应预见到氧气的存在；而供给氧气（或本可移除却未移除氧气）的人，通常并不被期待预见到划火柴的仪式。

然而，法律应在必要成分与充分成分之间赋予多大权重？这个问题显然超出我们考察的范围，而且究竟谁有资格处理这一问题、我们的法律体系是否准备实施相关建议，都完全不清楚。不过我希望，凡着手考虑此类问题的人，都会发现本章的分析有所助益。下一章把必要性与充分性的各个方面结合起来，阐明一个更精细的概念：“实际原因”。

\section*{致谢}
\addcontentsline{toc}{section}{致谢}

我感谢 Sander Greenland 就流行病学文献中归因的处理以及我们的结果在实际流行病学研究中的潜在应用所提出的许多建议与讨论。Donald Michie 与 Jack Good 促使我把注意力从 PN 转向 PS 与 PNS。Clark Glymour 与 Patricia Cheng 帮助解开了因果力理论的一些奥秘，Michelle Pearl 提供了流行病学文献的有用指引。Blai Bonet 纠正了引理 9.2.7 与 9.2.8 早期版本中的遗漏，Jin Tian 则用紧的界把一切收束起来。
